\section{Conclusiones}
\label{sec:conclusiones}

\textit{[PENDIENTE: discusión final estructurada en cuatro ejes:
(i) si los resultados obtenidos por las CNN propias permiten
considerar el sistema como una herramienta de apoyo diagnóstico
clínicamente confiable; (ii) análisis de las clases problemáticas
y sus posibles causas (similitud fotométrica, desbalance,
variabilidad intra-clase); (iii) limitaciones del enfoque sin
\textit{transfer learning} y comparación cualitativa con
literatura que sí lo emplea; y (iv) trabajo futuro, incluyendo
\textit{transfer learning} desde redes preentrenadas,
segmentación previa de la región de interés, ampliación del
dataset y validación con datos clínicos externos.]}

\subsection*{Hallazgos del Modelo 2}

El Modelo~2 (DeepCNN-BN y DeepCNN-BN-GAP) permite cerrar tres
preguntas planteadas a la luz de los resultados del Modelo~1:

\begin{enumerate}
  \item \textbf{¿Es \textit{Batch Normalization} responsable
        del colapso de OkanNet?} \textbf{No.} Bajo el régimen
        adaptado (AdamW, lr=$5\times 10^{-4}$, bs=$64$,
        wd=$10^{-4}$), una red con $10$ capas convolucionales
        y BN tras cada una alcanza F1-macro de $0{,}5961$ en
        \textit{Testing} ---$+17{,}28$ puntos porcentuales
        sobre OkanNet---, con curvas de convergencia
        monótonas y sin oscilaciones. El colapso de OkanNet
        era atribuible a la combinación lr=$10^{-3}$ /
        bs=$32$, no a BN como técnica. La hipótesis sostenida
        en la Sección~\ref{sec:metricas-modelo} sobre el
        Modelo~1 queda confirmada.

  \item \textbf{¿Aporta el cabezal FC capacidad
        discriminativa real frente a GAP?} \textbf{Sí, en
        este dataset.} La variante DeepCNN-BN supera a
        DeepCNN-BN-GAP por $7{,}85$ puntos de F1-macro en
        \textit{Testing} ($0{,}5961$ vs $0{,}5176$), pese a
        compartir el \textit{backbone} convolucional. La
        brecha se concentra en las clases con mayor
        variabilidad intra-clase (\textit{meningioma\_tumor}:
        $+14{,}6$ pp). La hipótesis de que GAP rinde
        equivalentemente cuando el \textit{backbone} es
        suficientemente profundo \emph{no se sostiene} en
        este corpus: los $\sim 524$\,K parámetros adicionales
        del cabezal FC sí se traducen en mejor generalización.

  \item \textbf{¿Resuelve el aumento de capacidad el sesgo
        \textit{glioma}~$\to$~\textit{meningioma}?}
        \textbf{No.} El Modelo~2 reproduce, e incluso
        intensifica, el patrón Precision-Recall observado en
        el baseline: alta Precision ($1{,}0000$ en DeepCNN-BN)
        y bajísimo Recall ($0{,}1600$) sobre
        \textit{glioma\_tumor}. Esto descarta una explicación
        puramente arquitectónica del problema y refuerza la
        hipótesis de divergencia distribucional entre los
        \textit{splits} oficiales del dataset Sartaj Bhuvaji:
        los gliomas del \textit{Testing} parecen provenir de
        subtipos o protocolos de adquisición sub-representados
        en \textit{Training}, lo cual el modelo no puede
        compensar sin intervención en los datos.
\end{enumerate}

Como corolario práctico, el Modelo~2 \textbf{no supera al
baseline en \textit{Testing}} ($0{,}5961$ vs $0{,}6252$,
$-2{,}91$ pp), pese a triplicar la profundidad convolucional y
añadir múltiples capas de regularización. Esta caída,
consistente con la pérdida de $\sim 30$ puntos entre validación
y prueba que afecta a todas las arquitecturas evaluadas en
este trabajo, indica que el factor dominante de generalización
\textbf{no es la capacidad del modelo sino el \textit{distribution
shift} entre particiones}. Direcciones esperables para el
Modelo~3 incluyen estrategias específicamente orientadas a este
problema, como \textit{transfer learning} desde redes
preentrenadas (que ya capturan invariancias robustas frente al
\textit{shift}), aumento de datos basado en transformaciones
físicamente plausibles para MRI, o esquemas de
\textit{domain adaptation} entre \textit{splits}.
